\begin{figure}[htbp]
	\centering
	\subfloat[][\emph{traslazione lungo la linea d'azione}]{%
		\begin{tikzpicture}[>=Stealth, line width=0.8pt, scale=0.80]
			% linea d'azione (direzione di fb_k)
			\draw[dashed] (0,0) -- (4.5,2.25);
			% punto Pk
			\fill (1,0.5) circle (1.5pt) node[below right] {$P_k$};
			% forza in Pk
			\draw[->, thick] (1,0.5) -- (2.6,1.3)
			node[above left] {$\fb_k$};
			% punto Pk' più avanti sulla stessa retta
			\fill (3,1.5) circle (1.5pt) node[below right] {$P_k'$};
			% forza in Pk'
			\draw[->, thick] (3,1.5) -- (4.6,2.3)
			node[above left] {$\fb_k$};
	\end{tikzpicture}}
	\qquad
	\subfloat[][\emph{composizione di forze concorrenti}]{%
		\begin{tikzpicture}[>=Stealth, line width=0.8pt, scale=0.80]
			% punto di concorrenza P
			\fill (2,1.5) circle (1.5pt) node[above right] {$P$};
			% linee d'azione
			\draw[dashed] (0,0) -- (2,1.5);
			\draw[dashed] (4,0) -- (2,1.5);
			% forze f1 e f2 dirette verso P
			\draw[->, thick] (0.6,0.45) -- (2,1.5)
			node[midway, left] {$\fb_1$};
			\draw[->, thick] (3.4,0.45) -- (2,1.5)
			node[midway, right] {$\fb_2$};
			% risultante in P
			\draw[->, thick] (2,1.5) -- (2,3.0)
			node[right] {$\fb_1+\fb_2$};
	\end{tikzpicture}}
	\qquad
	\subfloat[][\emph{decomposizione di una forza}]{%
		\begin{tikzpicture}[>=Stealth, line width=0.8pt, scale=0.80]
			% punto P
			\fill (1.5,1.0) circle (1.5pt) node[below right] {$P$};
			% forza risultante f
			\draw[->, thick] (1.5,1.0) -- (3.3,2.5)
			node[above right] {$\fb$};
			% componente f1 (orizzontale)
			\draw[->, thick] (1.5,1.0) -- (3.3,1.0)
			node[below] {$\fb_1$};
			% componente f2 (verticale)
			\draw[->, thick] (1.5,1.0) -- (1.5,2.5)
			node[left] {$\fb_2$};
			% parallelogramma tratteggiato
			\draw[dashed] (3.3,1.0) -- (3.3,2.5);
			\draw[dashed] (1.5,2.5) -- (3.3,2.5);
	\end{tikzpicture}}
	\caption{Operazioni elementari di equivalenza (I).
		\textbf{\textit{(a)}}~Traslazione della forza $\fb_k$
		lungo la propria linea d'azione da $P_k$ a $P_k'$:
		risultante e momento risultante restano invariati.
		\textbf{\textit{(b)}}~Composizione di due forze
		concorrenti in $P$: il sistema è equivalente alla
		risultante $\fb_1+\fb_2$ applicata in $P$.
		\textbf{\textit{(c)}}~Decomposizione della forza $\fb$
		applicata in $P$ nelle componenti $\fb_1$ e $\fb_2$
		secondo due direzioni indipendenti: operazione inversa
		alla composizione.}
	\label{fig:operazioni_elementari_1}
\end{figure}
